\section{Discussion}
\label{sec:discussion}

This section separates observations supported by the development runs from
hypotheses that require human labels, M2 implementation, and repeated final
evaluation.

\subsection{Development Observations}

\paragraph{Syntax is a necessary but weak success criterion.}

The transport correction demonstrates a basic methodological point: reliable
structured output is necessary before scientific evaluation is meaningful.
Once the JSON boundary is stable, however, parseability says little about
schema compliance, target consistency, plausibility, redundancy, coverage, or
utility. Future comparisons should therefore report transport and parse
failures separately from quality findings rather than rewarding a generator
for syntax alone.

\paragraph{Record validity and dataset quality lead to different decisions.}

The pilot B2/B3 disagreements are partly explained by policy rather than by
critic accuracy. B2 rejects repeated canonical candidates; the blinded human
task explicitly does not mark an otherwise sound record invalid merely because
it is redundant. A record-validity benchmark must therefore remove
duplicate-only B2 rejections from the structural-validity view, while a dataset
construction benchmark should retain them as evidence of low marginal value.
This separation is central to \rcqasyn{} and will be tested directly after
human annotation.

\paragraph{Self-judging and example anchoring.}

B3's near-universal acceptance in P1 cannot yet be interpreted as either
correct tolerance or critic failure. The judge and generator share a model,
and the human labels required to distinguish these explanations are pending.
The reference-copying diagnostics raise a parallel concern: class-aware
examples may improve conditioning while anchoring the generator to a small
visible support set. The magnitude varies across the three pilot datasets,
which argues for reporting reference identity and copying rather than treating
in-context examples as an unexamined implementation detail.

\paragraph{Scale and redundancy.}

Duplicates increased between P0-v3 and the larger P1 request for all three
development datasets. With one run at each scale, the study cannot establish a
trend or explain its cause.

\subsection{Hypotheses for Later Stages}

No implemented signal answers every evaluation question. Deterministic rules
are precise within their declared domain but do not establish semantic
plausibility; a language-model critic is flexible but potentially correlated
and uncalibrated; copying and support diagnostics describe dataset contact but
do not define record validity; downstream utility aggregates many upstream
effects. We therefore hypothesize that scoped heterogeneous evidence will be
more informative than any single critic and that calibration will improve
decision reliability. We separately hypothesize that dataset-aware control can
improve coverage and that localized repair can reduce action cost. None of
these hypotheses is established by P1; the latter two may remain extensions if
they are unnecessary for the core result.

% TODO[DISCUSSION]: revisit every interpretive claim after human adjudication,
% repeated splits, final benchmarks, and M2+ ablations.
